\exercisesection{\S~6}

\begin{enumerate}
\def\labelenumi{\arabic{enumi})}
\item
  Show that \(Z(\lambda)\) can be defined as a quotient of the representation \(\rho_{\lambda}\) of §4, Exerc. 1.
\item
  Let \(\mu\) be a weight of \(Z(\lambda)\) (resp. \(E(\lambda)\) ). Show that there exists a sequence of weights \(\mu_0, \ldots, \mu_n\) of \(Z(\lambda)\) (resp. of \(E(\lambda)\) ) such that \(\mu_0 = \lambda, \mu_n = \mu\) , and \(\mu_{i-1} - \mu_i \in B\) for \(1 \leq i \leq n\) .
\item
  Assume that g is simple, and denote by \(\tilde{\alpha}\) the highest root of R. The module \(\mathrm{E}(\tilde{\alpha})\) is isomorphic to g, equipped with the adjoint representation. If C is the Casimir element associated to the Killing form of g, the image of C in \(\mathrm{End}\mathrm{E}(\tilde{\alpha})\) is the identity (cf.~Chap. I, §3, no. 7, Prop. 12). Deduce, by using the Cor. of Prop. 7, that
\end{enumerate}

\[
\varPhi_ {\mathrm{R}} (\tilde {\alpha}, \tilde {\alpha} + 2 \rho) = 1,
\]

where \(\Phi_{\mathrm{R}}\) is the canonical bilinear form on \(\mathfrak{h}^*\) (Chap. VI, §1, no. 12).

\begin{enumerate}
\def\labelenumi{\arabic{enumi})}
\setcounter{enumi}{3}
\tightlist
\item
  We use the notations of no. 4.
\end{enumerate}

\begin{enumerate}
\def\labelenumi{\alph{enumi})}
\item
  Let \(m \in \mathbf{N}\) . Give \(\mathbf{N}^m\) the product order. Show that, for any subset \(S\) of \(\mathbf{N}^m\) , the set of minimal elements of \(S\) is finite.
\item
  Let \(\alpha_{1},\ldots ,\alpha_{m}\) be distinct elements of R, \(X_{i}\in \mathfrak{g}^{\alpha_{i}} - \{0\}\) , S the set of non-zero sequences \((p_i)\in \mathbf{N}^m\) such that \(\sum p_i\alpha_i = 0\) , M the set of minimal elements of S. Then \(\mathfrak{h}\) and the \(X_1^{p_1}\dots X_m^{p_m}\) , where \((p_1,\dots ,p_m)\in \mathrm{M}\) , generate the algebra \(\mathrm{U}^0\) .
\item
  Show that \(U^{0}\) is both a left- and right-noetherian algebra. (Give \(U^{0}\) the filtration induced by that of \(\mathrm{U}(\mathfrak{g})\) , and show, by using a) and b), that gr \(U^{0}\) is commutative of finite type.)
\end{enumerate}

\(d\) ) Show that, for all \(\lambda \in \mathfrak{h}^*\) , \(\mathrm{U}^{\lambda}\) is a left (resp. right) \(\mathrm{U}^0\) -module of finite type.

\begin{enumerate}
\def\labelenumi{\alph{enumi})}
\setcounter{enumi}{4}
\tightlist
\item
  Let V be a simple g-module such that \(V = \bigoplus_{\lambda \in b^{*}} V_{\lambda}\) . If one of the \(V_{\lambda} \neq 0\) is finite dimensional, then all of the \(V_{\lambda}\) are finite dimensional. (Use d).
\end{enumerate}

\begin{enumerate}
\def\labelenumi{\arabic{enumi})}
\setcounter{enumi}{4}
\tightlist
  \item
  Show that, if \(\mathfrak{g} = \mathfrak{sl}(2,k)\) , the modules \(Z(\lambda)\) of this paragraph are isomorphic to the modules \(Z(\lambda)\) of §1, Exerc. 2.
\end{enumerate}

